% Standalone document that typesets the five figures of the numerical study
% (Section 4 of "Positivity of an SIR model with delay") from the tables in
% pics/data/.  Build from the repository root with
%
%     latexmk -pdf figures.tex        (or: pdflatex figures.tex)
%
% after regenerating the data with  python -m experiments.run_all  (run inside
% python/).  The figure sources pics/fig_*.tex are identical to the ones
% included by the manuscript.
\documentclass[a4paper]{article}
\usepackage{a4wide}
\usepackage{xcolor}
\usepackage{amssymb,amsmath}
\usepackage{tikz}
\usepackage{pgfplots}
\pgfplotsset{compat=1.18}
\usepgfplotslibrary{groupplots,fillbetween}

% Colours used by pics/sirstyle.tex (copied from the manuscript's def.tex).
\definecolor{lightgray}{RGB}{175,175,175}
\definecolor{mycolor0}{rgb}{0.66,0.66,0.66}% grey
\definecolor{mycolor4}{rgb}{0.00000,0.44700,0.74100}% blue
\definecolor{mycolor1}{rgb}{0.85000,0.32500,0.09800}% orange
\definecolor{mycolor2}{rgb}{0.92900,0.69400,0.12500}% yellow
\definecolor{mycolor3}{rgb}{0.67000,0.74700,0.14100}% green
\definecolor{mycolor5}{rgb}{0.49400,0.18400,0.55600}% violet

\newcommand{\calR}{\mathcal{R}}

% The legends cross-reference the manuscript; resolve those references to the
% numbers they carry there instead of loading cleveref.
\makeatletter
\newcommand{\Cref}[1]{\@ifundefined{sirref@#1}{#1}{\@nameuse{sirref@#1}}}
\@namedef{sirref@thm:positivity}{Theorem~2}
\@namedef{sirref@tab:params}{Table~1}
\makeatother

% Shared pgfplots styles for the numerical study.
% All figures read the tables written by python/experiments into pics/data/.
\pgfplotsset{
    sirbase/.style={
        width=0.37\textwidth,
        height=0.30\textwidth,
        scale only axis,
        tick label style={font=\scriptsize},
        label style={font=\small},
        title style={font=\small},
        legend style={font=\scriptsize, draw=none, fill=none,
                      cells={anchor=west}, inner sep=1pt},
        every axis plot/.append style={line width=0.8pt},
        grid=major,
        grid style={lightgray!40, line width=0.3pt},
        axis line style={lightgray!80!black},
    },
    sirS/.style={color=mycolor4},              % susceptible  (blue)
    sirI/.style={color=mycolor1},              % infectious   (orange)
    sirR/.style={color=mycolor3},              % removed      (green)
    sirN/.style={color=black},                 % total population
    sirTheory/.style={color=mycolor5, dashed}, % published bound
    sirSharp/.style={color=mycolor2, dashed, line width=1.3pt},
    sirData/.style={only marks, mark=*, mark size=0.6pt, color=mycolor0},
}

% Generated by the experiments in python/experiments -- do not edit.
% Experiment 0 -- verification.
\newcommand{\expZeroRateAcademic}{1.02}
\newcommand{\expZeroRateCovid}{1.02}
\newcommand{\expZeroRuns}{44}
\newcommand{\expZeroMinSIR}{1.07\cdot 10^{-20}}
\newcommand{\expZeroMaxN}{1.1111}
\newcommand{\expZeroBound}{1.1111}

% Generated by the experiments in python/experiments -- do not edit.
% Experiment 1 -- sharpness of the positivity conditions.
\newcommand{\expOneSmallTauRatioMax}{11.2}
\newcommand{\expOneSplitErrReal}{0.5}
\newcommand{\expOneSplitMaxReal}{24}
\newcommand{\expOneBetaCrit}{4.904}

% Generated by the experiments in python/experiments -- do not edit.
% Experiment 2 -- consequences of neglecting the delay.
\newcommand{\expTwoRZero}{2.00}
\newcommand{\expTwoRateZero}{0.099}
\newcommand{\expTwoRateThirty}{0.018}
\newcommand{\expTwoPeakShift}{528}
\newcommand{\expTwoPeakDrop}{64}
\newcommand{\expTwoFitRZero}{1.38}
\newcommand{\expTwoFitRZeroBias}{31}
\newcommand{\expTwoFitLogErr}{0.00}
\newcommand{\expTwoPeakPrevErr}{-58}
\newcommand{\expTwoAttackErr}{-43}
\newcommand{\expTwoCovidTauLo}{4.5}
\newcommand{\expTwoWinLo}{58}
\newcommand{\expTwoWinHi}{79}

% Generated by the experiments in python/experiments -- do not edit.
% Experiment 3 -- real-world data.
\newcommand{\expThreeRItaly}{0.200}
\newcommand{\expThreeRItalyLo}{0.189}
\newcommand{\expThreeRItalyHi}{0.220}
\newcommand{\expThreeRLiberia}{0.0475}
\newcommand{\expThreeCovidRode}{2.53}
\newcommand{\expThreeCovidRdelay}{6.89}
\newcommand{\expThreeCovidFactor}{2.72}
\newcommand{\expThreeEbolaRode}{1.48}
\newcommand{\expThreeEbolaRdelay}{2.37}
\newcommand{\expThreeBetaDelay}{0.90}
\newcommand{\expThreeRtable}{0.105}
\newcommand{\expThreeLockDay}{14}
\newcommand{\expThreeTzeroOde}{25}
\newcommand{\expThreeTzeroDelay}{21}
\newcommand{\expThreeAbove}{3.4}
\newcommand{\expThreeBelow}{2.8}
\newcommand{\expThreeBias}{-0.05}
\newcommand{\expThreeLockRedOde}{70}
\newcommand{\expThreeLockRedDelay}{91}
\newcommand{\expThreeLockBetaOne}{1.41}
\newcommand{\expThreeLockBetaTwo}{0.13}
\newcommand{\expThreeLockSse}{15.1}
\newcommand{\expThreeLockPeak}{18}
\newcommand{\expThreeDataPeak}{28}
\newcommand{\expThreeLagBest}{3}
\newcommand{\expThreeLagSse}{12.7}
\newcommand{\expThreeRmseOde}{0.29}
\newcommand{\expThreeRmseDelay}{0.35}
\newcommand{\expThreeRedOde}{67}
\newcommand{\expThreeRedDelay}{90}
\newcommand{\expThreeBetaOdeOne}{0.35}
\newcommand{\expThreeBetaDelayOne}{0.98}

% Generated by the experiments in python/experiments -- do not edit.
% Experiment 4 -- tau -> 0 limit and regularity.
\newcommand{\expFourTauRate}{1.05}
\newcommand{\expFourJumpOne}{1.3\cdot 10^{-6}}
\newcommand{\expFourJumpTwo}{2.2\cdot 10^{-2}}


\begin{document}

\begin{figure}[p]
	\centering
	% Experiment 0.2 and 4a: convergence in the step size and in the delay.
\begin{tikzpicture}
\begin{groupplot}[
    group style={group size=2 by 1, horizontal sep=2.0cm},
    sirbase, xmode=log, ymode=log,
]
\nextgroupplot[
    xlabel={step size $h$}, 
    ylabel={error at $t=T$},
    xmin = 6e-4,
    xmax = 4,
    ymin = 2e-6,
    ymax = 0.098,
    legend pos=north west,
]
\addplot[color=mycolor4, mark=*, mark size=2.2pt]
    table[x=h, y=error] {pics/data/exp0_convergence_academic.dat};
\addlegendentry{sub-critical set, $\tau=1$}
\addplot[color=mycolor1, mark=square*, mark size=2.2pt]
    table[x=h, y=error] {pics/data/exp0_convergence_covid.dat};
\addlegendentry{COVID-19 set, $\tau=5$}
\addplot[sirTheory, domain=1e-3:1, samples=2, forget plot] {0.01*x};
\node[font=\scriptsize, anchor=west, mycolor5] at (axis cs:1,1e-2) {$\mathcal{O}(h)$};

\nextgroupplot[
    xlabel={delay $\tau$}, 
    %ylabel={distance to delay-free model},
    ylabel={$\|u_\tau(T)-u_0(T)\|_\infty$},
    legend pos=north west,
]
\addplot[color=mycolor4, mark=triangle*, mark size=2.8pt]
    table[x=tau, y=error] {pics/data/exp4a_taulimit.dat};
%\addlegendentry{$\|u_\tau(T)-u_0(T)\|_\infty$}
\addplot[sirTheory, domain=3.5e-3:0.1, samples=2, forget plot] {2e-3*x};
\node[font=\scriptsize, anchor=south east, mycolor5] at (axis cs:0.1,1.8e-4)
    {$\mathcal{O}(\tau)$};
\end{groupplot}
\end{tikzpicture}

	\caption{(Manuscript Figure~1.) Left: first-order convergence, error at
	$t=T$ against the step size~$h$. Right: distance to the solution of the
	delay-free model as $\tau\to0$.}
\end{figure}

\begin{figure}[p]
	\centering
	% Experiment 1b/1d: the true positivity boundary against the sufficient conditions.
% The measured curves are drawn first so that the two reference curves stay
% visible where the measured boundary settles exactly on top of them.
\begin{tikzpicture}
\begin{groupplot}[
    group style={group size=2 by 1, horizontal sep=2.0cm},
    sirbase, xmode=log, ymode=log,
]
\nextgroupplot[
    xlabel={transmission rate $\beta$},
    ylabel={critical history level $\bar N_0$},
    legend pos=north east,
    legend style={fill=white, fill opacity=0.85, text opacity=1},
]
\addplot[color=mycolor0]  table[x=beta, y=crit0p1] {pics/data/exp1b_boundary.dat};
\addlegendentry{measured, $\tau=0.1$}
\addplot[color=mycolor3]  table[x=beta, y=crit1]   {pics/data/exp1b_boundary.dat};
\addlegendentry{measured, $\tau=1$}
\addplot[color=mycolor4]  table[x=beta, y=crit10]  {pics/data/exp1b_boundary.dat};
\addlegendentry{measured, $\tau=10$}
%\addplot[color=mycolor4, densely dashdotted]
%    table[x=beta, y=crit100] {pics/data/exp1b_boundary.dat};
%\addlegendentry{measured, $\tau=100$}
\addplot[sirSharp]   table[x=beta, y=thm]   {pics/data/exp1b_boundary.dat};
\addlegendentry{$2\sqrt{\Lambda/\beta}$ (\Cref{thm:positivity})}

\nextgroupplot[
    xlabel={delay $\tau$}, 
    ylabel={critical history level $\bar N_0$},
    ymode=normal,
    xmin=0.1,
    xmax=100,
    ymin=1.8, 
    ymax=8, 
    legend pos=north east,
    legend style={fill=white, fill opacity=0.85, text opacity=1},
]
\addplot[color=mycolor1, mark=*, mark size=1.5pt]
    table[x=tau, y=crit]  {pics/data/exp1d_tau.dat};
\addlegendentry{measured, $\beta=0.8$}
\addplot[sirSharp]  table[x=tau, y=thm] {pics/data/exp1d_tau.dat};
\addlegendentry{$2\sqrt{\Lambda/\beta}$ (\Cref{thm:positivity})}
\end{groupplot}
\end{tikzpicture}

	\caption{(Manuscript Figure~2.) Sharpness of the history condition. Left:
	measured critical history level $\bar N_0$ against $\beta$ for three
	delays. Right: the same threshold as a function of $\tau$ at $\beta=0.8$.}
\end{figure}

\begin{figure}[p]
	\centering
	% Experiment 2a/2b: identical R_0, different delay; and the calibration mismatch.
\begin{tikzpicture}
\begin{groupplot}[
    group style={group size=2 by 1, horizontal sep=2.0cm},
    sirbase,
]
\nextgroupplot[
    xlabel={time $t$ [days]}, ylabel={infectious $I(t)$},
    xmin=0, xmax=800, legend pos=north east,
]
\addplot[color=mycolor4] table[x=t, y=I0]  {pics/data/exp2a_delays.dat};
\addlegendentry{$\tau=0$}
\addplot[color=mycolor3] table[x=t, y=I5]  {pics/data/exp2a_delays.dat};
\addlegendentry{$\tau=5$}
\addplot[color=mycolor1] table[x=t, y=I10] {pics/data/exp2a_delays.dat};
\addlegendentry{$\tau=10$}
\addplot[color=mycolor5] table[x=t, y=I20] {pics/data/exp2a_delays.dat};
\addlegendentry{$\tau=20$}
\addplot[color=mycolor2] table[x=t, y=I30] {pics/data/exp2a_delays.dat};
\addlegendentry{$\tau=30$}

\nextgroupplot[
    xlabel={time $t$ [days]}, ylabel={infectious $I(t)$}, ymode=log,
    xmin=0, xmax=800, ymin=5e-7, ymax=0.95, legend pos=north west,
    legend style={fill=white, fill opacity=0.85, text opacity=1},
]
% calibration window, shaded
\addplot[draw=none, fill=mycolor0!25, forget plot] coordinates
    {(\expTwoWinLo,1e-7) (\expTwoWinHi,1e-7) (\expTwoWinHi,0.3) (\expTwoWinLo,0.3)}
    \closedcycle;
\addplot[color=mycolor1] table[x=t, y=Itrue]  {pics/data/exp2b_mismatch.dat};
\addlegendentry{delay model, truth}
\addplot[color=mycolor4, densely dashed]
    table[x=t, y=Inaive] {pics/data/exp2b_mismatch.dat};
\addlegendentry{delay-free fit}
\node[font=\scriptsize, anchor=south west, align=left]
    at (axis cs:105,1e-6) {calibration\\ window};
\end{groupplot}
\end{tikzpicture}

	\caption{(Manuscript Figure~3.) Effect of the delay at fixed
	$\calR_0 = \expTwoRZero$. Left: prevalence for several delays. Right: a
	delay-free model calibrated on the shaded window to data generated with
	$\tau=10$.}
\end{figure}

\begin{figure}[p]
	\centering
	% Experiment 2c/2d: loss of positivity at realistic parameters, and delay-induced
% oscillations inside the sharpened admissible region.
\begin{tikzpicture}
\begin{groupplot}[
    group style={group size=2 by 1, horizontal sep=2.0cm},
    sirbase,
]
\nextgroupplot[
    xlabel={delay $\tau$ [days]}, ylabel={$\min_t S(t)$},
    xmin=0, xmax=40, legend pos=north west,
    legend style={fill=white, fill opacity=0.85, text opacity=1},
]
\addplot[color=mycolor3, line width=1.5pt] table[x=tau, y=minS_ebola] {pics/data/exp2c_positivity.dat};
\addlegendentry{Ebola set (\Cref{tab:params})}
\addplot[color=mycolor1, line width=1.5pt] table[x=tau, y=minS_covid] {pics/data/exp2c_positivity.dat};
\addlegendentry{COVID-19 set (\Cref{tab:params})}
\addplot[black, line width=0.4pt, forget plot] coordinates {(0,0) (40,0)};

\nextgroupplot[
    xlabel={time $t$}, ylabel={infectious $I(t)$},
    xmin=300, xmax=400, ymin=0.51, ymax=0.76,
    legend pos=north east, legend columns=4,
    legend style={fill=white, fill opacity=0.85, text opacity=1,
                  /tikz/every even column/.append style={column sep=4pt}},
]
\addplot[color=mycolor4, line width=1.4pt] table[x=t, y=I0]
    {pics/data/exp2d_hopf.dat};
\addlegendentry{$\tau=0$\ }
\addplot[color=mycolor3, line width=1.4pt] table[x=t, y=I3] {pics/data/exp2d_hopf.dat};
\addlegendentry{$\tau=3$\ }
\addplot[color=mycolor1, line width=1.4pt] table[x=t, y=I4] {pics/data/exp2d_hopf.dat};
\addlegendentry{$\tau=4$}
%\addplot[color=mycolor5] table[x=t, y=I7] {pics/data/exp2d_hopf.dat};
%\addlegendentry{$\tau=7$}
\end{groupplot}
\end{tikzpicture}

	\caption{(Manuscript Figure~4.) Left: minimum of $S$ against the delay for
	the two parameter sets of Table~1. Right: delay-induced oscillations for
	the super-critical dimensionless set.}
\end{figure}

\begin{figure}[p]
	\centering
	% Experiment 3c: the piecewise-constant beta(t) fits to the Italian incidence.
% Thick curves: model incidence; thin curves with marks: the same centred 7-day
% mean as for the data, applied to the model incidence.
\begin{tikzpicture}
\begin{groupplot}[
    %group style={group size=2 by 1, horizontal sep=1.5cm},
    sirbase,
    width=0.7\textwidth,
]
\nextgroupplot[
    xlabel={days after 24 February 2020}, ylabel={daily incidence [density]},
    ymode=log, xmin=0, xmax=70, ymin=8e-7, ymax=6e-4,
    legend columns=2,
    legend style={at={(0.5,-0.2)}, anchor=north,
                  /tikz/every even column/.append style={column sep=10pt}},
]
\addplot[sirSharp, forget plot] coordinates {(14,8e-7) (14,6e-4)};
\addplot[sirData, mark size=1.4pt]
    table[x=day, y=data] {pics/data/exp3c_fits.dat};
\addlegendentry{reported (7-day mean)}
\addlegendimage{empty legend}
\addlegendentry{}
\addplot[color=mycolor4, line width=1.5pt, densely dashed]
    table[x=day, y=ode] {pics/data/exp3c_fits.dat};
\addlegendentry{no delay, $t_0=\expThreeTzeroOde$}
\addplot[color=mycolor4, line width=0.6pt, mark=square*, mark size=1pt]
    table[x=day, y=ode_avg] {pics/data/exp3c_fits.dat};
\addlegendentry{no delay, 7-day mean}
\addplot[color=mycolor1, line width=1.5pt]
    table[x=day, y=delay] {pics/data/exp3c_fits.dat};
\addlegendentry{delay, fitted $t_0$}
\addplot[color=mycolor1, line width=0.6pt, mark=*, mark size=1pt]
    table[x=day, y=delay_avg] {pics/data/exp3c_fits.dat};
\addlegendentry{delay, fitted $t_0$, 7-day mean}
\addplot[color=mycolor3, line width=1.5pt, densely dashdotted]
    table[x=day, y=lock] {pics/data/exp3c_fits.dat};
\addlegendentry{delay, $t_0$ = lockdown}
\addplot[color=mycolor3, line width=0.6pt, mark=triangle*, mark size=1pt]
    table[x=day, y=lock_avg] {pics/data/exp3c_fits.dat};
\addlegendentry{delay, $t_0$ = lockdown, 7-day mean}
\node[font=\scriptsize, mycolor2, anchor=south west, rotate=90] at (axis cs:17.3,9e-7)
    {lockdown, 9 Mar};

\end{groupplot}
\end{tikzpicture}

	\caption{(Manuscript Figure~5.) Reported Italian daily incidence
	(seven-day mean) and piecewise-constant $\beta(t-\tau)$ fits with a fitted
	change point and with the change point pinned to the lockdown.}
\end{figure}

\end{document}
